\subsection{Модель реалистичного рендеринга}

В общем случае задача компьютерной графики заключается в отображении трехмерных данных на двумерном сенсоре (мониторе), при этом спект решаемых таким образом задач очень сильно варируется. Например, компьютерная графика находит большое применение в автоматизированном проектировании (CAD системы).

Область, которой касается настоящая работа, является проблема реалистичного рендеринга, то есть построения системы, симулирующей поведение света.

Современный рендеринг разделяет поведение света на два случая:
\begin{itemize}
	\item На границе раздела двух сред с разным коэффициентом приломления часть света отражается, часть приломляется; математической моделью этого процесса в рендеринге выступают двунаправленные функции распределения рассеивания (BSDF), за отражение и приломление отвечают соотвественно двунаправленные функции отражения и пропускания (BRDF, BTDF);
	\item В участвующих средах свет непосредственно взаимодействует и подвергается поглощению и рассеиванию, это описывается уравнением переноса излучения (RTE), о котором пойдет речь в разделе \ref{sec:rte}.
\end{itemize}

Наиболее популярными движками рендера являются RenderMan и Arlond (используются для офлайн рендеринга в кинематографе), Unreal Engine и Unity (рендеринг в реальном времени), Cycles (решение с открытым исходным кодом). Отдельно стоит выделить движок $Mitsuba$, реализующий дифференцируемый рендеринг.